\begin{filecontents*}{refs.bib}
@article{ho2020ddpm,
  title={Denoising Diffusion Probabilistic Models},
  author={Ho, Jonathan and Jain, Ajay and Abbeel, Pieter},
  journal={NeurIPS},
  year={2020}
}
@article{song2020score,
  title={Score-Based Generative Modeling through Stochastic Differential Equations},
  author={Song, Yang and Sohl-Dickstein, Jascha and Kingma, Diederik P. and Kumar, Abhishek and Ermon, Stefano and Poole, Ben},
  journal={ICLR},
  year={2021}
}
@article{song2020ddim,
  title={Denoising Diffusion Implicit Models},
  author={Song, Jiaming and Meng, Chenlin and Ermon, Stefano},
  journal={ICLR},
  year={2021}
}
@article{nichol2021improved,
  title={Improved Denoising Diffusion Probabilistic Models},
  author={Nichol, Alexander Quinn and Dhariwal, Prafulla},
  journal={ICML},
  year={2021}
}
@article{karras2022edm,
  title={Elucidating the Design Space of Diffusion-Based Generative Models},
  author={Karras, Tero and Aittala, Miika and Aila, Timo and Laine, Samuli},
  journal={NeurIPS},
  year={2022}
}
@article{kingma2021variational,
  title={Variational Diffusion Models},
  author={Kingma, Diederik P. and Salimans, Tim and Poole, Ben and Ho, Jonathan},
  journal={NeurIPS},
  year={2021}
}
@article{dhariwal2021guided,
  title={Diffusion Models Beat GANs on Image Synthesis},
  author={Dhariwal, Prafulla and Nichol, Alexander Quinn},
  journal={NeurIPS},
  year={2021}
}
@article{peebles2023dit,
  title={Scalable Diffusion Models with Transformers},
  author={Peebles, William and Xie, Saining},
  journal={ICCV},
  year={2023}
}
@article{saharia2022imagen,
  title={Photorealistic Text-to-Image Diffusion Models with Deep Language Understanding},
  author={Saharia, Chitwan and Chan, William and Saxena, Saurabh and Li, Lala and Whang, Jay and Denton, Emily and others},
  journal={NeurIPS},
  year={2022}
}
@article{rombach2022ldm,
  title={High-Resolution Image Synthesis with Latent Diffusion Models},
  author={Rombach, Robin and Blattmann, Andreas and Lorenz, Dominik and Esser, Patrick and Ommer, Bjorn},
  journal={CVPR},
  year={2022}
}
\end{filecontents*}

\documentclass{article}
\usepackage[preprint]{neurips_2024}
\usepackage[utf8]{inputenc}
\usepackage[T1]{fontenc}
\usepackage{hyperref}
\usepackage{url}
\usepackage{booktabs}
\usepackage{amsfonts}
\usepackage{nicefrac}
\usepackage{microtype}
\usepackage{xcolor}
\usepackage{amsmath,amssymb}
\usepackage{graphicx}

\newcommand{\placeholder}[1]{\textcolor{red}{#1}}

\title{Adaptive Noise Scheduling for Class-Conditional Diffusion Models}
\author{Anonymous Authors}
\date{}

\begin{document}
\maketitle

\begin{abstract}
Diffusion models typically use a fixed noise schedule shared across all samples, despite large variation in semantic complexity across classes. We introduce \textbf{Adaptive Noise Scheduling} (ANS), a class-conditional framework that learns $\beta_t(y)$ using a dedicated ScheduleNet and jointly optimizes generation fidelity, schedule efficiency, and temporal smoothness. ScheduleNet receives class embedding and timestep encoding, then predicts bounded per-step noise levels with differentiable constraints enforcing valid diffusion dynamics. Our objective augments standard $\epsilon$-prediction loss with (i) an efficiency term that discourages unnecessary cumulative noise and (ii) a smoothness regularizer stabilizing schedule trajectories. For evaluation, we use a fair paired-checkpoint protocol: adaptive and fixed-cosine models are trained separately under matched architecture and optimization settings. We report FID--time frontiers with repeated sampling runs and uncertainty estimates, emphasizing the correct claim: adaptive scheduling should improve quality-time tradeoffs by reaching comparable fidelity at fewer denoising steps, not by reducing per-step compute at identical step counts.
\end{abstract}

\section{Introduction}
Diffusion models have become a dominant paradigm for high-fidelity generation~\cite{ho2020ddpm,nichol2021improved,dhariwal2021guided}. A central but under-adapted component is the forward noise schedule $\{\beta_t\}_{t=1}^T$, usually fixed (linear or cosine) for all data points. This global schedule ignores heterogeneity in data complexity: classes with simple structure may not require the same perturbation profile as classes with fine texture and shape ambiguity.

We propose \textbf{Adaptive Noise Scheduling (ANS)}, where a neural scheduler outputs class-specific $\beta_t(y)$ and is trained jointly with the diffusion denoiser.

Contributions:
\begin{itemize}
  \item A differentiable ScheduleNet that predicts bounded class-conditional noise schedules.
  \item A multi-objective training formulation combining diffusion fidelity with efficiency and smoothness constraints.
  \item A practical evaluation protocol reporting FID, runtime, and per-class adaptive-vs-fixed schedule behavior.
\end{itemize}

\section{Background}
Standard DDPM forward diffusion defines:
\begin{align}
q(\mathbf{x}_t \mid \mathbf{x}_{t-1}) &= \mathcal{N}\!\left(\sqrt{1-\beta_t}\mathbf{x}_{t-1}, \beta_t \mathbf{I}\right), \\
\bar{\alpha}_t &= \prod_{i=1}^{t} (1-\beta_i), \\
\mathbf{x}_t &= \sqrt{\bar{\alpha}_t}\mathbf{x}_0 + \sqrt{1-\bar{\alpha}_t}\,\boldsymbol{\epsilon}, \quad \boldsymbol{\epsilon}\sim\mathcal{N}(0,\mathbf{I}).
\end{align}
Training minimizes
\begin{align}
\mathcal{L}_{\text{diffusion}} = \mathbb{E}_{t,\mathbf{x}_0,\boldsymbol{\epsilon}}\left[\left\|\boldsymbol{\epsilon} - \boldsymbol{\epsilon}_\theta(\mathbf{x}_t,t)\right\|_2^2\right].
\end{align}

\section{Method}
\subsection{Adaptive Schedule Parameterization}
Let $g_\phi$ denote ScheduleNet. Given class embedding $\mathbf{c}(y)$ and timestep $t$, we predict
\begin{align}
\beta_t(y) = g_\phi\!\left(\mathbf{c}(y), \text{PE}(t/T)\right), \quad \beta_t(y)\in[\beta_{\min},\beta_{\max}],
\end{align}
with learnable bounds satisfying $0<\beta_{\min}<\beta_{\max}\le 0.02$.
Then
\begin{align}
\bar{\alpha}_t(y) = \prod_{i=1}^{t}\left(1-\beta_i(y)\right).
\end{align}

\subsection{Joint Objective}
We optimize:
\begin{align}
\mathcal{L} &= \mathcal{L}_{\text{diffusion}}
 + \lambda_1 \mathcal{L}_{\text{efficiency}}
 + \lambda_2 \mathcal{L}_{\text{smoothness}},
\end{align}
where
\begin{align}
\mathcal{L}_{\text{efficiency}} &= \mathbb{E}_{y}\left[\sum_{t=1}^T \beta_t(y)\right], \\
\mathcal{L}_{\text{smoothness}} &= \mathbb{E}_{y}\left[\sum_{t=1}^{T-1}\left(\beta_{t+1}(y)-\beta_t(y)\right)^2\right].
\end{align}

\paragraph{Why $\mathcal{L}_{\text{efficiency}}$ can improve the frontier.}
For fixed denoiser capacity, large cumulative noise implies larger reverse correction magnitude. Minimizing $\sum_t \beta_t(y)$ reduces average corruption budget and can preserve sample fidelity under coarser reverse discretizations (fewer DDIM steps). We therefore evaluate the method through quality-time frontier shifts rather than same-step speed comparisons.

\section{Experiments}
\subsection{Setup}
Dataset: CIFAR-10 ($32\times32$). Base denoiser: class-conditional U-Net. ScheduleNet: residual MLP with sinusoidal timestep encoding. We train \emph{two separate checkpoints}: (i) adaptive schedule mode and (ii) fixed cosine schedule mode, using matched architecture and optimization settings.

\subsection{Baselines}
Primary baseline is a separately trained fixed-cosine model~\cite{nichol2021improved}. We additionally report linear schedules and DDIM step sweeps~\cite{song2020ddim}. We do \emph{not} treat fixed-sampler counterfactuals from adaptive-trained weights as primary baselines.

\subsection{Metrics}
Primary metrics: FID (with large-sample protocol), wall-clock sampling time, per-class FID/time, and schedule diversity (pairwise class schedule distance). For all key tables we report mean$\pm$std over repeated sampling runs.

\section{Results}
\begin{table}[t]
\centering
\caption{Main quality/efficiency comparison on CIFAR-10 (separately trained checkpoints, mean$\pm$std over repeated runs).}
\begin{tabular}{lcc}
\toprule
Method & FID $\downarrow$ & Time (s / 64 imgs) $\downarrow$ \\
\midrule
Fixed cosine DDIM-50 & \placeholder{xx.xx} & \placeholder{x.xxx} \\
Adaptive DDIM-50 & \placeholder{xx.xx} & \placeholder{x.xxx} \\
Fixed DDPM-1000 & \placeholder{xx.xx} & \placeholder{x.xxx} \\
Adaptive DDPM-1000 & \placeholder{xx.xx} & \placeholder{x.xxx} \\
\bottomrule
\end{tabular}
\end{table}

\begin{table}[t]
\centering
\caption{Per-class adaptive-vs-fixed summary (class-matched real references, mean$\pm$std).}
\begin{tabular}{lccc}
\toprule
Class & FID$_{\text{adaptive}}$ & FID$_{\text{fixed}}$ & Speedup \\
\midrule
airplane & \placeholder{} & \placeholder{} & \placeholder{} \\
automobile & \placeholder{} & \placeholder{} & \placeholder{} \\
bird & \placeholder{} & \placeholder{} & \placeholder{} \\
cat & \placeholder{} & \placeholder{} & \placeholder{} \\
deer & \placeholder{} & \placeholder{} & \placeholder{} \\
dog & \placeholder{} & \placeholder{} & \placeholder{} \\
frog & \placeholder{} & \placeholder{} & \placeholder{} \\
horse & \placeholder{} & \placeholder{} & \placeholder{} \\
ship & \placeholder{} & \placeholder{} & \placeholder{} \\
truck & \placeholder{} & \placeholder{} & \placeholder{} \\
\bottomrule
\end{tabular}
\end{table}

Figure~\ref{fig:frontier} presents the efficiency frontier; Figure~\ref{fig:perclass} shows per-class speedup.
\begin{figure}[t]
  \centering
  \includegraphics[width=0.75\linewidth]{figures/efficiency_frontier.png}
  \caption{Adaptive vs fixed FID-time frontier across denoising step counts.}
  \label{fig:frontier}
\end{figure}
\begin{figure}[t]
  \centering
  \includegraphics[width=0.75\linewidth]{figures/per_class_speedup.png}
  \caption{Per-class speedup ratio for adaptive scheduling.}
  \label{fig:perclass}
\end{figure}

\section{Ablation Study}
We ablate:
\begin{itemize}
  \item $\lambda_1$ (efficiency): $\{0, 10^{-3}, 10^{-2}, 10^{-1}\}$.
  \item $\lambda_2$ (smoothness): $\{0, 10^{-3}, 10^{-2}, 10^{-1}\}$.
  \item ScheduleNet depth: $\{1,2,3,4\}$ residual layers.
\end{itemize}
We report effects on FID, runtime, schedule diversity, and training stability.

\section{Conclusion}
Adaptive class-conditional noise scheduling is a simple, differentiable extension of DDPM training that targets the fidelity-efficiency tradeoff. With fair paired-checkpoint evaluation, ANS should be interpreted as a method for improving the FID-time frontier (comparable quality at fewer steps), rather than as a same-step runtime acceleration mechanism.

\bibliographystyle{plain}
\bibliography{refs}

\end{document}
